\chapter*{Prólogo}
\addcontentsline{toc}{chapter}{Prólogo}
\label{ch:prologo}

Durante la formación en la licenciatura en física, el estudio analítico de la estructura atómica suele culminar con la solución exacta del átomo de hidrógeno (o iones hidrogenoides), ya que la resolución analítica de la ecuación de Schrödinger está estrictamente limitada a sistemas de un solo electrón. Fuera de los cursos de especialidad, el problema de muchos cuerpos se aborda de manera introductoria; a menudo se menciona que al escalar hacia átomos más pesados resulta indispensable implementar correcciones relativistas, pero rara vez se detalla formalmente ni el \textit{cómo} ni el \textit{por qué} emergen estas interacciones.

Para sistemas multielectrónicos reales, la interacción repulsiva interelectrónica acopla las ecuaciones de movimiento, imposibilitando una solución exacta de forma cerrada. Esta intrincada correlación dinámica y la necesidad de ir más allá del modelo hidrogenoide han motivado el desarrollo histórico de la química cuántica computacional y la física atómica teórica, cuyo pilar fundamental es el método de Campo Autoconsistente de Hartree-Fock (HF).

El método de Hartree-Fock busca la mejor aproximación uniconfiguracional a la función de onda exacta utilizando un determinante de Slater, reduciendo el problema de muchos cuerpos a un sistema de ecuaciones integro-diferenciales no lineales. Tradicionalmente, para resolver estas ecuaciones, los orbitales atómicos se expanden en conjuntos de bases analíticas, como los orbitales de tipo Slater (STO) o funciones Gaussianas (GTO). Si bien estas bases globales son eficientes, sufren de limitaciones críticas: la sobrecompletitud, que genera problemas de dependencia lineal, y una representación deficiente en las regiones muy cercanas al núcleo y en las colas asintóticas de la función de onda.

%%% Local Variables:
%%% mode: LaTeX
%%% TeX-master: "../main"
%%% End:
